\subsection{The degree-six theorem and the remaining cases}
\begin{theorem}
Every singular modulus of degree at least six requires at least five nonzero coefficients in every nonzero rational polynomial that vanishes at it. For degree at least four, the same conclusion holds except possibly at the following 79 absolute discriminants. The assertion also holds over the corresponding CM field.
\end{theorem}
\begin{center}\small
\begin{tabular}{@{}c p{.84\linewidth}@{}}
\toprule Degree & Remaining absolute discriminants \\ \midrule
4 & 39, 55, 56, 63, 68, 80, 128, 136, 144, 155, 156, 171, 184, 196, 203, 208, 219, 220, 252, 256, 259, 275, 291, 292, 323, 328, 355, 363, 387, 388, 400, 475, 507, 568, 592, 603, 667, 723, 763, 772, 955, 1003, 1027, 1227, 1243, 1387, 1411, 1467, 1507, 1555. \\
5 & 47, 79, 103, 127, 131, 179, 188, 227, 316, 347, 412, 443, 508, 523, 571, 619, 683, 691, 739, 787, 947, 1051, 1123, 1723, 1747, 1867, 2203, 2347, 2683. \\
\bottomrule\end{tabular}\end{center}
\begin{proof}
The analytic cutoff handles $D\ge10^{11}$, and the intermediate sieve covers $10^7\le D<10^{11}$. The small sieve gives precisely the list above together with $4075$ in degree at least four. The preceding complex determinant computation removes that final degree-six case. Linear disjointness gives the assertion over the CM field. The candidate list records where the proofs stop. It does not assert the existence of a four-term relation at any listed discriminant.
\end{proof}
